% GENERATED FILE. Edit canonical lessons, then run npm run generate:books.
% canonical-source-hash: fnv1a64:e6bba9dcc6ae5218
% canonical-lessons: BN-C437-bagh, BN-C437-horin, BN-C437-shikar-kora, BN-C437-chora, BN-C437-gungun-kora, BN-R437-review-431-437, BN-R437-review-final-437

\chapter{Wild Creatures}
\label{ch:bn-wild-creatures}
\hlchaptermodality{437}

\begin{chapteropening}
\textbf{By the end of this chapter:} \emph{I can talk about tigers and deer, hunting and grazing, and the hum of insects.}

Complete the last lesson of chapter 437: I can talk about tigers and deer, hunting and grazing, and the hum of insects.
\end{chapteropening}

\section[\texorpdfstring{\bn{বাঘ} (bāgh)}{বাঘ (bāgh)}]{\bn{বাঘ} (bāgh) — a tiger}

\label{lesson:BN-C437-bagh}



\begin{quote}
Before the new one: say the Bengali for to sprout, then the Bengali for to plant (in the ground).
\end{quote}

\subsection*{You'll want to know: \bn{বাঘ}}
\textbf{\bn{বাঘ}} — \emph{bāgh} — \textquotedblleft{}a tiger\textquotedblright{}.

\begin{quote}
\bn{জঙ্গলে} \bn{আমরা} \bn{একটা} \bn{বাঘ} \bn{দেখলাম।}
\end{quote}

\begin{quote}
\emph{jônggôle āmrā ekṭā bāgh dekhlām.}
\end{quote}

\begin{quote}
We saw a tiger in the forest.
\end{quote}

The tiger of the Sundarbans mangroves is the pride of Bengal.

\begin{cousinweb}
\textbf{\bn{বাঘ}} (\emph{bāgh}) is inherited: CDIAL lists Bengali \emph{bāgh}, \textquotedblleft{}a tiger\textquotedblright{}, under Sanskrit \emph{vyāghrá}, \textquotedblleft{}a tiger\textquotedblright{} (CDIAL 12193).
\end{cousinweb}

\begin{grammarlens}[title={What you've built}]
One more word for getting around.
\end{grammarlens}

\subsection*{Your turn}
\begin{itemize}
\raggedright
  \item \emph{Say it:} \emph{bāgh}
  \item \emph{Say it:} \emph{bāgh}, once more
  \item \emph{Say it:} \emph{jônggôle āmrā ekṭā bāgh dekhlām.} — We saw a tiger in the forest.
  \item \emph{Recall:} say the Bengali for to sprout, then the Bengali for to plant (in the ground), then say \emph{bāgh} again
\end{itemize}

\begin{tcolorbox}[breakable,colback=teal!4,colframe=teal!35!black,title={Before you move on}]
What does \bn{বাঘ} mean? (\textquotedblleft{}A tiger\textquotedblright{}.) Say it once more.
\end{tcolorbox}
\section[\texorpdfstring{\bn{হরিণ} (hôrin)}{হরিণ (hôrin)}]{\bn{হরিণ} (hôrin) — a deer}

\label{lesson:BN-C437-horin}



\begin{quote}
Before the new one: say the Bengali for to plant (in the ground), then the Bengali for a tiger.
\end{quote}

\subsection*{You'll want to know: \bn{হরিণ}}
\textbf{\bn{হরিণ}} — \emph{hôrin} — \textquotedblleft{}a deer\textquotedblright{}.

\begin{quote}
\bn{বাঘ} \bn{দেখে} \bn{হরিণগুলো} \bn{দৌড়ে} \bn{পালাল।}
\end{quote}

\begin{quote}
\emph{bāgh dekhe hôringulo douṛe pālālô.}
\end{quote}

\begin{quote}
Seeing the tiger, the deer ran away.
\end{quote}

\begin{cousinweb}
\textbf{\bn{হরিণ}} (\emph{hôrin}) is inherited: CDIAL lists Bengali \emph{harin}, \textquotedblleft{}a deer\textquotedblright{}, under Sanskrit \emph{hariṇá}, \textquotedblleft{}yellowish; a deer\textquotedblright{} (CDIAL 13982).
\end{cousinweb}

\begin{grammarlens}[title={What you've built}]
One more word for getting around.
\end{grammarlens}

\subsection*{Your turn}
\begin{itemize}
\raggedright
  \item \emph{Say it:} \emph{hôrin}
  \item \emph{Say it:} \emph{hôrin}, once more
  \item \emph{Say it:} \emph{bāgh dekhe hôringulo douṛe pālālô.} — Seeing the tiger, the deer ran away.
  \item \emph{Recall:} say the Bengali for to plant (in the ground), then the Bengali for a tiger, then say \emph{hôrin} again
\end{itemize}

\begin{tcolorbox}[breakable,colback=teal!4,colframe=teal!35!black,title={Before you move on}]
What does \bn{হরিণ} mean? (\textquotedblleft{}A deer\textquotedblright{}.) Say it once more.
\end{tcolorbox}
\section[\texorpdfstring{\bn{শিকার} \bn{করা} (shikār kôrā)}{শিকার করা (shikār kôrā)}]{\bn{শিকার} \bn{করা} (shikār kôrā) — to hunt}

\label{lesson:BN-C437-shikar-kora}



\begin{quote}
Before the new one: say the Bengali for a tiger, then the Bengali for a deer.
\end{quote}

\subsection*{You'll want to know: \bn{শিকার} \bn{করা}}
\textbf{\bn{শিকার} \bn{করা}} — \emph{shikār kôrā} — \textquotedblleft{}to hunt\textquotedblright{}.

\begin{quote}
\bn{বাঘ} \bn{রাতে} \bn{শিকার} \bn{করে।}
\end{quote}

\begin{quote}
\emph{bāgh rāte shikār kôre.}
\end{quote}

\begin{quote}
Tigers hunt at night.
\end{quote}

Past (he/she …): \bn{সেই} \bn{বাঘ} \bn{একটা} \bn{হরিণ} \bn{শিকার} \bn{করল}, \emph{sei bāgh ekṭā hôrin shikār kôrlô}.

\bn{শিকার} alone is the hunt, and also the prey.

\begin{cousinweb}
\textbf{\bn{শিকার} \bn{করা}} (\emph{shikār kôrā}) joins \bn{শিকার}, usually traced to Persian \emph{shikār}, \textquotedblleft{}hunting, the game hunted\textquotedblright{}, to \bn{করা}, to do.
\end{cousinweb}

\begin{grammarlens}[title={What you've built}]
One more word for getting around.
\end{grammarlens}

\subsection*{Your turn}
\begin{itemize}
\raggedright
  \item \emph{Say it:} \emph{shikār kôrā}
  \item \emph{Say it:} \emph{shikār kôrā}, once more
  \item \emph{Say it:} \emph{bāgh rāte shikār kôre.} — Tigers hunt at night.
  \item \emph{Recall:} say the Bengali for a tiger, then the Bengali for a deer, then say \emph{shikār kôrā} again
\end{itemize}

\begin{tcolorbox}[breakable,colback=teal!4,colframe=teal!35!black,title={Before you move on}]
What does \bn{শিকার} \bn{করা} mean? (\textquotedblleft{}To hunt\textquotedblright{}.) Say it once more.
\end{tcolorbox}
\section[\texorpdfstring{\bn{চরা} (chôrā)}{চরা (chôrā)}]{\bn{চরা} (chôrā) — to graze}

\label{lesson:BN-C437-chora}



\begin{quote}
Before the new one: say the Bengali for a deer, then the Bengali for to hunt.
\end{quote}

\subsection*{You'll want to know: \bn{চরা}}
\textbf{\bn{চরা}} — \emph{chôrā} — \textquotedblleft{}to graze\textquotedblright{}.

\begin{quote}
\bn{মাঠে} \bn{গরু} \bn{চরছিল।}
\end{quote}

\begin{quote}
\emph{māṭhe gôru chôrchhilô.}
\end{quote}

\begin{quote}
Cows were grazing in the field.
\end{quote}

Past (he/she …): \bn{সেই} \bn{গরু} \bn{মাঠে} \bn{চরল}, \emph{sei gôru māṭhe chôrlô}.

\bn{চরানো}, to take cattle out to graze, is the causing partner.

\begin{cousinweb}
\textbf{\bn{চরা}} (\emph{chôrā}) is inherited: CDIAL lists Bengali \emph{carā}, \textquotedblleft{}to move about, to graze\textquotedblright{}, under Sanskrit \emph{cárati}, \textquotedblleft{}moves, grazes\textquotedblright{} (CDIAL 4686).
\end{cousinweb}

\begin{grammarlens}[title={What you've built}]
One more word for getting around.
\end{grammarlens}

\subsection*{Your turn}
\begin{itemize}
\raggedright
  \item \emph{Say it:} \emph{chôrā}
  \item \emph{Say it:} \emph{chôrā}, once more
  \item \emph{Say it:} \emph{māṭhe gôru chôrchhilô.} — Cows were grazing in the field.
  \item \emph{Recall:} say the Bengali for a deer, then the Bengali for to hunt, then say \emph{chôrā} again
\end{itemize}

\begin{tcolorbox}[breakable,colback=teal!4,colframe=teal!35!black,title={Before you move on}]
What does \bn{চরা} mean? (\textquotedblleft{}To graze\textquotedblright{}.) Say it once more.
\end{tcolorbox}
\section[\texorpdfstring{\bn{গুনগুন} \bn{করা} (gungun kôrā)}{গুনগুন করা (gungun kôrā)}]{\bn{গুনগুন} \bn{করা} (gungun kôrā) — to hum; to buzz}

\label{lesson:BN-C437-gungun-kora}



\begin{quote}
Before the new one: say the Bengali for to hunt, then the Bengali for to graze.
\end{quote}

\subsection*{You'll want to know: \bn{গুনগুন} \bn{করা}}
\textbf{\bn{গুনগুন} \bn{করা}} — \emph{gungun kôrā} — \textquotedblleft{}to hum; to buzz\textquotedblright{}.

\begin{quote}
\bn{রাতে} \bn{কানের} \bn{কাছে} \bn{মশা} \bn{গুনগুন} \bn{করছিল।}
\end{quote}

\begin{quote}
\emph{rāte kāner kāchhe môshā gungun kôrchhilô.}
\end{quote}

\begin{quote}
At night a mosquito was buzzing by my ear.
\end{quote}

Past (he/she …): \bn{তিনি} \bn{গান} \bn{গুনগুন} \bn{করলেন}, \emph{tini gān gungun kôrlen}.

Of insects, and of a person humming a tune under the breath: \bn{গান} \bn{গুনগুন} \bn{করা}.

\begin{cousinweb}
\textbf{\bn{গুনগুন} \bn{করা}} (\emph{gungun kôrā}) joins \bn{গুনগুন}, an imitative word for a low hum or buzz, to \bn{করা}, to do. No dictionary we use records an older form, so its history is not settled.
\end{cousinweb}

\begin{grammarlens}[title={What you've built}]
One more word for getting around.
\end{grammarlens}

\subsection*{Your turn}
\begin{itemize}
\raggedright
  \item \emph{Say it:} \emph{gungun kôrā}
  \item \emph{Say it:} \emph{gungun kôrā}, once more
  \item \emph{Say it:} \emph{rāte kāner kāchhe môshā gungun kôrchhilô.} — At night a mosquito was buzzing by my ear.
  \item \emph{Recall:} say the Bengali for to hunt, then the Bengali for to graze, then say \emph{gungun kôrā} again
\end{itemize}

\begin{tcolorbox}[breakable,colback=teal!4,colframe=teal!35!black,title={Before you move on}]
What does \bn{গুনগুন} \bn{করা} mean? (\textquotedblleft{}To hum; to buzz\textquotedblright{}.) Say it once more.
\end{tcolorbox}
\section[(dialogue)]{Looking back: from Bays, Creeks and Deltas to Wild Creatures}

\label{lesson:BN-R437-review-431-437}



\begin{quote}
Say the Bengali for a bay and a creek.
\end{quote}

\subsection*{Your turn}
Say these aloud:

\begin{itemize}
\raggedright
  \item a bay, a creek, a delta, silt, to fill in (a pond or low land)
  \item a summit, a valley, a mound, a desert, a jungle
  \item a lake (natural), a ditch, a wetland, to sway, to babble (of water)
  \item nature, the environment, pollution, to spoil, to conserve
  \item grass, a creeper, a branch, wood, to trim
  \item a sprout, a mango blossom, to bloom, to sprout, to plant (in the ground)
  \item a tiger, a deer, to hunt, to graze, to hum
\end{itemize}

\begin{tcolorbox}[breakable,colback=teal!4,colframe=teal!35!black,title={Before you move on}]
A bay? (\textbf{\bn{উপসাগর}}, \emph{upôsāgôr}.) A creek? (\textbf{\bn{খাঁড়ি}}, \emph{khā̃ṛi}.) A delta? (\textbf{\bn{বদ্বীপ}}, \emph{bôdvīp}.) Silt? (\textbf{\bn{পলি}}, \emph{pôli}.) To fill in (a pond or low land)? (\textbf{\bn{ভরাট} \bn{করা}}, \emph{bhôrāṭ kôrā}.) A summit? (\textbf{\bn{চূড়া}}, \emph{chūṛā}.) A valley? (\textbf{\bn{উপত্যকা}}, \emph{upôtyôkā}.) A mound? (\textbf{\bn{ঢিবি}}, \emph{ḍhibi}.) A desert? (\textbf{\bn{মরুভূমি}}, \emph{môrubhūmi}.) A jungle? (\textbf{\bn{জঙ্গল}}, \emph{jônggôl}.) A lake (natural)? (\textbf{\bn{ঝিল}}, \emph{jhil}.) A ditch? (\textbf{\bn{নালা}}, \emph{nālā}.) A wetland? (\textbf{\bn{জলাভূমি}}, \emph{jôlābhūmi}.) To sway? (\textbf{\bn{দোল} \bn{খাওয়া}}, \emph{dol khāoyā}.) To babble (of water)? (\textbf{\bn{কলকল} \bn{করা}}, \emph{kôlkôl kôrā}.) Nature? (\textbf{\bn{প্রকৃতি}}, \emph{prôkriti}.) The environment? (\textbf{\bn{পরিবেশ}}, \emph{pôribesh}.) Pollution? (\textbf{\bn{দূষণ}}, \emph{dūshôn}.) To spoil? (\textbf{\bn{নষ্ট} \bn{করা}}, \emph{nôshṭô kôrā}.) To conserve? (\textbf{\bn{সংরক্ষণ} \bn{করা}}, \emph{sôngrôkkhôn kôrā}.) Grass? (\textbf{\bn{ঘাস}}, \emph{ghās}.) A creeper? (\textbf{\bn{লতা}}, \emph{lôtā}.) A branch? (\textbf{\bn{শাখা}}, \emph{shākhā}.) Wood? (\textbf{\bn{কাঠ}}, \emph{kāṭh}.) To trim? (\textbf{\bn{ছাঁটা}}, \emph{chhā̃ṭā}.) A sprout? (\textbf{\bn{অঙ্কুর}}, \emph{ôngkur}.) A mango blossom? (\textbf{\bn{মুকুল}}, \emph{mukul}.) To bloom? (\textbf{\bn{ফোটা}}, \emph{phoṭā}.) To sprout? (\textbf{\bn{গজানো}}, \emph{gôjāno}.) To plant (in the ground)? (\textbf{\bn{পোঁতা}}, \emph{põtā}.) A tiger? (\textbf{\bn{বাঘ}}, \emph{bāgh}.) A deer? (\textbf{\bn{হরিণ}}, \emph{hôrin}.) To hunt? (\textbf{\bn{শিকার} \bn{করা}}, \emph{shikār kôrā}.) To graze? (\textbf{\bn{চরা}}, \emph{chôrā}.) To hum? (\textbf{\bn{গুনগুন} \bn{করা}}, \emph{gungun kôrā}.)
\end{tcolorbox}
\section[(dialogue)]{Wild Creatures — once more}

\label{lesson:BN-R437-review-final-437}



\begin{quote}
Say the Bengali for a tiger and a deer.
\end{quote}

\subsection*{Your turn}
\begin{itemize}
\raggedright
  \item \emph{Say it:} a tiger, a deer, to hunt, to graze, to hum
\end{itemize}

\begin{tcolorbox}[breakable,colback=teal!4,colframe=teal!35!black,title={Before you move on}]
A tiger? (\textbf{\bn{বাঘ}}, \emph{bāgh}.) A deer? (\textbf{\bn{হরিণ}}, \emph{hôrin}.) To hunt? (\textbf{\bn{শিকার} \bn{করা}}, \emph{shikār kôrā}.) To graze? (\textbf{\bn{চরা}}, \emph{chôrā}.) To hum? (\textbf{\bn{গুনগুন} \bn{করা}}, \emph{gungun kôrā}.)
\end{tcolorbox}
